\documentclass[DIV=calc]{scrartcl}
\usepackage[utf8]{inputenc}
\usepackage[T1]{fontenc}
\usepackage[ngerman]{babel}
\usepackage{graphicx}
\usepackage[draft, markup=underlined]{changes}
\usepackage{csquotes}
\usepackage{eurosym}

\usepackage{ulem}
%\usepackage[dvipsnames]{xcolor}
\usepackage{paralist}
\usepackage{fixltx2e}
%\usepackage{ellipsis}
\usepackage[tracking=true]{microtype}

\usepackage{lmodern}              % Ersatz fuer Computer Modern-Schriften
%\usepackage{hfoldsty}

%\usepackage{fourier}             % Schriftart
\usepackage[scaled=0.81]{helvet}     % Schriftart

\usepackage{url}
%\usepackage{tocloft}             % Paket für Table of Contents
\def\UrlBreaks{\do\a\do\b\do\c\do\d\do\e\do\f\do\g\do\h\do\i\do\j\do\k\do\l%
\do\m\do\n\do\o\do\p\do\q\do\r\do\s\do\t\do\u\do\v\do\w\do\x\do\y\do\z\do\0%
\do\1\do\2\do\3\do\4\do\5\do\6\do\7\do\8\do\9\do\-}%

\usepackage{xcolor}
\definecolor{urlred}{HTML}{660000}

\usepackage{hyperref}
\hypersetup{colorlinks=false}

%\usepackage{mdwlist}     % Änderung der Zeilenabstände bei itemize und enumerate
% \usepackage{draftwatermark} % Wasserzeichen ``Entwurf''
% \SetWatermarkText{Antrag}

\parindent 0pt                 % Absatzeinrücken verhindern
\parskip 12pt                 % Absätze durch Lücke trennen

\setlength{\textheight}{23cm}
\usepackage{fancyhdr}
\pagestyle{fancy}
\fancyhead{} % clear all header fields
\cfoot{}
\lfoot{Zusammenkunft aller Physik-Fachschaften}
\rfoot{www.zapfev.de\\stapf@zapf.in}
\renewcommand{\headrulewidth}{0pt}
\renewcommand{\footrulewidth}{0.1pt}
\newcommand{\gen}{*innen}
\addto{\captionsngerman}{\renewcommand{\refname}{Quellen}}

%%%% Mit-TeXen Kommandoset
\usepackage[normalem]{ulem}
\usepackage{xcolor}
\usepackage{xspace} 

\newcommand{\replace}[2]{
    \sout{\textcolor{blue}{#1}}~\textcolor{blue}{#2}}
\newcommand{\delete}[1]{
    \sout{\textcolor{red}{#1}}}
\newcommand{\add}[1]{
    \textcolor{blue}{#1}}

\newif\ifcomments
\commentsfalse
%\commentstrue

\newcommand{\red}[1]{{\ifcomments\color{red} {#1}\else{#1}\fi}\xspace}
\newcommand{\blue}[1]{{\ifcomments\color{blue} {#1}\else{#1}\fi}\xspace}
\newcommand{\green}[1]{{\ifcomments\color{green} {#1}\else{#1}\fi}\xspace}

\newcommand{\repl}[2]{{\ifcomments{\color{red} \sout{#1}}{\color{blue} {\xspace #2}}\else{#2}\fi}}
%\newcommand{\repl}[2]{{\color{red} \sout{#1}\xspace{\color{blue} {#2}}\else{#2}\fi}\xspace}

\newcommand{\initcomment}[2]{%
	\expandafter\newcommand\csname#1\endcsname{%
		\def\thiscommentname{#1}%
		\definecolor{col}{rgb}{#2}%
		\def\thiscommentcolor{col}%
}}

% initcomment Name RGB-color
\initcomment{Philipp}{0, 0.5, 0}

%\renewcommand{\comment}[1]{{\ifcomments{\color{red} {#1}}{}\fi}\xspace}

\renewcommand{\comment}[2][\nobody]{
	\ifdefined#1
	{\ifcomments{#1 \expandafter\color{\thiscommentcolor}{\thiscommentname: #2}}{}\fi}\xspace
	\else
	{\ifcomments{\color{red} {#2}}{}\fi}\xspace
	\fi
}

\newcommand{\zapf}{ZaPF\xspace}

\let\oldgrqq=\grqq
\def\grqq{\oldgrqq\xspace}

\setlength{\parskip}{.6em}
\setlength{\parindent}{0mm}

%\usepackage{geometry}
%\geometry{left=2.5cm, right=2.5cm, top=2.5cm, bottom=3.5cm}

% \renewcommand{\familydefault}{\sfdefault}




\begin{document}

\hspace{0.87\textwidth}
\begin{minipage}{120pt}
	\vspace{-1.8cm}
	\includegraphics[width=80pt]{../logo.pdf}
	\centering
	\small Zusammenkunft aller Physik-Fachschaften
\end{minipage}

\begin{center}
  \huge{Resolution BAföG - aktuelle Überarbeitung}\vspace{.25\baselineskip}\\
  \normalsize
\end{center}
\vspace{1cm}

%%%% Metadaten %%%%

%\paragraph{Adressierte:} Alle im dt. Bundestag vertretenen Parteien und deren bildungspolitischen Sprecher:innen, Familienministerium, BMBF


%\paragraph{Antragstellende:} Christian (Marburg), Timo (Freiburg)

%%%% Text des Antrages zur veröffentlichung %%%%

%\section*{Antragstext}

Im Koalitionsvertrag\footnote{\url{https://www.bundesregierung.de/breg-de/service/gesetzesvorhaben/koalitionsvertrag-2021-1990800}} der Ampelparteien heißt es zum Thema BAföG: \\
\glqq Mit (...) einem grundlegend reformierten BAföG legen wir den Grundstein für ein Jahrzehnt der Bildungschancen.\grqq
Weiter heißt es:
\glqq Wir wollen allen Menschen unabhängig von ihrer Herkunft beste Bildungschancen bieten, Teilhabe und Aufstieg ermöglichen und durch inklusive Bildung sichern. Dazu (...) machen [wir] das BAföG elternunabhängiger und bauen es für die Förderung der beruflichen Weiterbildung aus\grqq

Wir begrüßen die Anerkennung dessen, dass beim BAföG Reformbedarf besteht. Leider wird auch die geplante Änderung der komplexen Lebensrealitäten der Studierenden nicht gerecht. Der aktuelle Entwurf zur 27. Novelle des Bundesausbildungsgesetzes\footnote{\url{https://dserver.bundestag.de/btd/20/017/2001734.pdf}} wird nicht zu einem grundlegend reformierten BAföG führen. 

Dennoch wollen wir zwei Punkte, die deutliche Verbesserungen mit sich bringen, positiv hervorheben:

Die ZaPF sieht die geplante Anhebung der Altersgrenze, für die Berechtigung auf BAföG, auf 45 Jahre sehr positiv. Ebenfalls spricht sich die ZaPF für die geplante Anhebung der Vermögensfreigrenze der Antragssteller:innen auf 45.000 Euro aus. Bei beiden Punkten sieht die ZaPF schon lange Änderungsbedarf\footnote{\label{bafögreso}\url{https://zapfev.de/resolutionen/wise20/bafoeg/bafoeg.pdf}}.

Trotz dieser Verbesserungen kann die kommende Änderung des BAföGs nur als eine kleine Anpassung, anstelle der angekündigten, grundlegenden Änderungen gesehen werden. 

Die Erhöhung der Freibeträge auf das Elterneinkommen um 20\% gleicht die Versäumnisse der letzten Jahre aus, kann aber nicht als eine grundlegende Reform verstanden werden. Wir bemängeln die starke Kopplung an das Einkommen der Eltern und sehen hier keine wirkliche Verbesserung.

Die Erhöhung des Bedarfssatzes um 5\% reicht nicht einmal aus, um der aktuellen Inflationsrate entgegenzuwirken. Das BAföG sollte die Lebenshaltungskosten von Studierenden decken und ihnen eine Teilhabe am sozialen und kulturellen Leben ermöglichen. Dies ist aktuell nicht der Fall und wird sich auch durch die angekündigten Anpassungen nicht ändern. Die ZaPF spricht sich für eine automatische Anpassung der BAföG-Sätze unter Berücksichtigung von Faktoren wie Inflation oder Lebenshaltungskosten aus, wie es beispielsweise bereits bei Bundestagsdiäten der Fall ist.

 Von einem \glqq Grundstein für ein Jahrzehnt der Bildungschancen\grqq kann nicht die Rede sein, wenn Das BAföG nicht zum Leben reicht und der Großteil der Studierenden auch weiterhin nicht dazu berechtigt ist, BAföG zu beziehen.

Die ZaPF fordert sich von der Regierungskoalition eine echte Reform des BAföGs und verweist an dieser Stelle auf ihren Forderungskatalog\footnote{\label{bafögreso}\url{https://zapfev.de/resolutionen/wise20/bafoeg/bafoeg.pdf}}.

Ebefalls schließen wir uns den Forderung der Hochschulrektorenkonferenz (HRK)\footnote{\url{https://www.hrk.de/positionen/beschluss/detail/anforderungen-an-eine-weiterentwicklung-des-bundesausbildungsfoerderungsgesetzes-bafoeg/}} und des BAföG50\footnote{\url{https://bafoeg50.de/reform/}} an. Diese entsprechen der nötigen Bedürfnisse der in der Realtät vorkommenden Lebensumstände.



%\newpage
%\section*{Begründung}


\vspace{1cm} 

\vfill
\begin{flushright}
	Verabschiedet am 07. Juni 2022 \\
	auf der ZaPF in Bochum.
\end{flushright}

\end{document}
